\section{Finite accuracy and a covariance determinant}\label{sec:finite-quadratic}

The rank law describes a fixed system as its tolerance tends to zero.
A different question keeps the tolerances finite and asks for constants
uniform over the coefficients and the number of outputs. Covariance
geometry answers this question, including when the output tolerances
are unequal. Throughout this section
\[
 f_j(x)=\tfrac12x^TH_jx+a_j^Tx+b_j,\qquad x\in[0,1]^n,
 \qquad H_j=H_j^T.
\]
Let $\omega_n$ be the volume of the Euclidean unit ball, and set
\begin{equation}\label{eq:finite-constants}
 c_n=\max\{4,n/4\},\quad
 A_n=\log_2\!\bigl[\omega_n(n+2)^{n/2}\bigr]
                   +\tfrac n2\log_2c_n,\quad
 B_n=n\log_2 n+n.
\end{equation}
The Gaussian integral gives
$\omega_n\le(2\pi e/n)^{n/2}$, so $A_n+B_n=O(n\log(n+1))$.

\begin{theorem}[Finite covariance law]\label{thm:finite-covariance}
For positive tolerances $\eps_j$, define
\begin{equation}\label{eq:finite-determinant}
 D_* =\max\left\{\det P:0\preceq P\preceq I,
   \ E_j(P):=\tr(H_jPH_jP)\le\eps_j^2\ (1\le j\le m)\right\},
 \qquad \Phi=-\tfrac12\log_2D_*.
\end{equation}
Then $0<D_*\le1$ and
\begin{equation}\label{eq:finite-law}
 \max\{0,\Phi-A_n\}\le\pconv\le\pbin\le\Phi+B_n.
\end{equation}
The upper construction has a binary count $p_{\rm grid}\le\Phi+B_n$
and $O(np_{\rm grid}+n^2+mn^2)$ rows and continuous variables.
The statement permits real coefficients; a rational polynomial-time
construction is proved in \cref{thm:rational-finite}.
\end{theorem}
\begin{proof}
The feasible set in \eqref{eq:finite-determinant} is compact and contains
a positive scalar multiple of $I$, so a determinant maximizer is positive
definite. Let $S$ be a positive-volume parity contact from
\cref{lem:parity}, with covariance $\Sigma$. The midpoint identity and
\cref{lem:covariance} give
$E_j(\Sigma)\le16\eps_j^2$. A unit linear functional on the cube has
range at most $\sqrt n$, hence variance at most $n/4$; thus
$\Sigma\preceq(n/4)I$. Consequently $\Sigma/c_n$ is feasible and
$\det\Sigma\le c_n^nD_*$. The volume inequality
\eqref{eq:covariance-volume} bounds every contact by
$2^{A_n}\sqrt{D_*}$. The same bound holds for zero-volume contacts.
The at most $2^p$ contacts cover a cube of volume one, proving the lower
bound without any measurability assumption on the initial lift projection.

For the upper bound, take any feasible
$P=U\diag(\lambda_i)U^T\succ0$ with $U$ orthogonal. After rotation and
translation, the cube lies in $0\le y_i\le w_i$, where
$w_i=\sum_k|U_{ki}|\in[1,\sqrt n]$. Keep the original cube inequalities
and the affine coordinate identities. The transformed Hessians are
$G_j=U^TH_jU$. Choose
\begin{equation}\label{eq:covariance-grid}
 L_i=\left\lceil\log_2\!\left(w_i\sqrt{n/\lambda_i}\right)\right\rceil,
 \qquad h_i=w_i2^{-L_i}\le\sqrt{\lambda_i/n}.
\end{equation}
Use the prefix and residual construction \eqref{eq:prefix}, scaled by
$w_i$. The exact binary products and residual envelopes
\eqref{eq:binary-product}--\eqref{eq:product-split} provide a shared
symmetric monomial approximation with errors $\delta_{ik}$ satisfying
$|\delta_{ik}|\le h_ih_k/4$. This includes squares, empty prefixes,
and the upper endpoints of the bounding intervals. The same approximate
monomial is used in every output. Thus
\[
 |w_j-f_j(x)|\le\tfrac18\sum_{i,k}|G_{j,ik}|h_ih_k
 \le\tfrac1{8n}\sum_{i,k}|G_{j,ik}|\sqrt{\lambda_i\lambda_k}
 \le\tfrac18\sqrt{E_j(P)}\le\eps_j/8.
\]
The last inequality before the energy uses Cauchy--Schwarz on $n^2$
entries. Choosing every approximate monomial equal to its exact value
proves graph containment. Finally,
\[
 \sum_i L_i\le-\tfrac12\log_2\det P
          +\sum_i\log_2w_i+\tfrac n2\log_2n+n
 \le-\tfrac12\log_2\det P+B_n.
\]
Use an optimal $P$. Each bit occurs in $O(n)$ binary products; there
are $O(n^2)$ residual monomials and $O(mn^2)$ output coefficients.
\end{proof}

On a full-dimensional compact convex $\Omega\subset[0,1]^n$, exactly
the same lower proof gives the useful variant
\begin{equation}\label{eq:finite-domain-volume}
 \pconv(F,\Omega;K)\ge\Phi-A_n+\log_2\vol(\Omega).
\end{equation}
Here $K$ is the tolerance box, and $\Phi$ uses the same Hessians on the
containing cube. The upper grid can be restricted to $\Omega$ by any
available exact linear lift of that domain. The volume term cannot in
general be discarded.

A feasible covariance also describes a directional-error ellipsoid. For
$d=P^{1/2}u$ with $\|u\|_2\le1$,
$|d^TH_jd|\le\|P^{1/2}H_jP^{1/2}\|_F\le\eps_j$.
Using the operator norm in place of the Frobenius norm changes the
log-determinant benchmark by at most $(n/4)\log_2n$: the norm inequalities
$\|M\|_2\le\|M\|_F\le\sqrt n\|M\|_2$ permit rescaling a
covariance by $1/\sqrt n$. This connects the finite invariant to
anisotropic quadratic interpolation and established derivative-based
ellipsoid metrics \citep{cao2007}, while the formulation comparison
still requires the contact and shared-grid arguments above.

\subsection{Geodesic convexity and finite certificates}

The feasible covariances need not be convex in the ordinary matrix
coordinates: a product Hessian on positive diagonal matrices has energy
$2p_1p_2$. They are convex along the positive definite geodesics
\[
 P(t)=P_0^{1/2}\exp(tA)P_0^{1/2},\qquad A=A^T.
\]
If $A=V\diag(a_i)V^T$ and $C_j=V^TP_0^{1/2}H_jP_0^{1/2}V$, then
\begin{equation}\label{eq:energy-geodesic}
 E_j(P(t))=\sum_{i,k}C_{j,ik}^2e^{t(a_i+a_k)}.
\end{equation}
Each energy is convex, and its logarithm is convex unless the Hessian
vanishes. The log determinant is affine. Between two endpoints the
geodesic obeys $P(t)\preceq(1-t)P(0)+tP(1)$: conjugate by $P(0)^{-1/2}$
and use $d^t\le1-t+td$ on each positive eigenvalue. Thus the cap and
all energy constraints are geodesically convex. Every local determinant
maximum in their positive definite part is global. These are applications
of established positive definite geometry; see \citet{sra2013}.

A feasible covariance is already an upper-construction certificate.
The following residual certificate also bounds its suboptimality and can
be used without executing the conservative algorithm below.
\begin{proposition}[A posteriori determinant certificate]\label{prop:covariance-certificate}
Let $0\prec P\preceq I$ satisfy every energy budget. For $\mu_j\ge0$
and $S\succeq0$, put
\[
 c=\sum_j\mu_j(\eps_j^2-E_j(P))+\tr[S(I-P)],\qquad
 R=-P^{-1}+2\sum_j\mu_jH_jPH_j+S,\qquad L=-\log\det P.
\]
Then $c\ge0$, and, with natural logarithms,
\begin{equation}\label{eq:covariance-certificate}
 0\le\log D_* -\log\det P
 \le c+\sqrt n\,L\sqrt{\tr(PRPR)}.
\end{equation}
In particular, $c=0$ and $R=0$ certify global optimality. For rational
$P,\mu,S$, the residual square is an exactly computable nonnegative
rational number; scalar logarithm and square-root intervals certify the
right side.
\end{proposition}
\begin{proof}
The Lagrangian
$J(X)=-\log\det X+\sum_j\mu_j(E_j(X)-\eps_j^2)+\tr[S(X-I)]$
is geodesically convex by \eqref{eq:energy-geodesic} and the corresponding
positive-coefficient expansion of $\tr(SX)$. Its tangent gradient at $P$,
in isometric coordinates, is $G=P^{1/2}RP^{1/2}$. An optimal feasible
$P_*$ has $\det P_*\ge\det P$. Since both matrices are capped by $I$,
their minimum eigenvalues are at least $\det P$. Therefore the eigenvalues
of $P^{-1/2}P_*P^{-1/2}$ lie between $\det P$ and $(\det P)^{-1}$.
Writing $A=\log(P^{-1/2}P_*P^{-1/2})$, we obtain
$\|A\|_F\le\sqrt nL$. Convexity gives
\[
 -\log D_*\ge J(P_*)\ge J(P)+\tr(GA)
 \ge-\log\det P-c-\sqrt nL\|G\|_F.
\]
Now $\|G\|_F^2=\tr(PRPR)$, proving the claim. For fixed $P$, the
right side is convex in the multipliers, since its slack is linear and
its residual norm is a norm of an affine map. No necessity of stationarity
or constraint qualification is asserted.
\end{proof}

\subsection{Exact reductions and a dimension gap}

For an edge output $x_ix_k$, direct multiplication gives
$E_{ik}(P)=2(P_{ii}P_{kk}+P_{ik}^2)$. Replacing $P$ by its diagonal
preserves the cap, decreases these energies and increases its determinant
by Hadamard's inequality. With $P_{ii}=2^{-2s_i}$, the benchmark becomes
\begin{equation}\label{eq:covariance-graph-lp}
 \Phi=\min\left\{\sum_i s_i:s_i\ge0,\quad
 s_i+s_k\ge\log_2(\sqrt2/\eps_{ik})\text{ for every edge}\right\}.
\end{equation}
This recovers the fractional-cover allocation of \cref{thm:graph-cover}
with its finite tolerance constants.

\begin{proposition}[Commuting Hessians]\label{prop:commuting-covariance}
If all Hessians commute, the benchmark has a maximizer diagonal in a
common orthonormal eigenbasis. Writing its eigenvalues as $p_i$ and the
Hessian eigenvalues as $h_{ji}$, it reduces to
\[
 \max\left\{\sum_i\log t_i:0<t_i\le1,
                   \ \sum_i h_{ji}^2t_i\le\eps_j^2\right\},
 \qquad t_i=p_i^2,
\]
whose optimal objective is $2\log D_*$. Hessians may be indefinite.
\end{proposition}
\begin{proof}
In the common eigenbasis every coordinate sign flip $S$ commutes with
every Hessian. Thus $P$ and $SPS$ have the same determinant and energies.
Their geometric mean $M=P\#(SPS)$ is feasible by geodesic convexity and
has determinant $\det P$. The identity $SMS=M$ follows from symmetry
and orthogonal congruence invariance of the geometric mean. These
identities follow in turn from its unique positive definite solution of
$MP^{-1}M=Q$: this equation also implies $MQ^{-1}M=P$.
A previously imposed commuting sign symmetry is preserved. Apply the
operation to the $n$ coordinate flips in succession; the result is
an optimal diagonal covariance. Substitution proves the scalar program.
\end{proof}

For one Hessian put $a_i=h_i^2$. If $\eps^2\ge\sum_i a_i$, then
$D_*=1$. Otherwise the unique $s>0$ solving
$\sum_i\min(a_i,s)=\eps^2$ gives
$t_i=\min(1,s/a_i)$ for $a_i>0$ and $t_i=1$ for $a_i=0$.
Indeed, the budget multiplier $1/s$ and upper-cap multipliers
$\max(0,1-a_i/s)$ satisfy the stationarity equations for the strictly
concave scalar objective. Sorting the $a_i$ determines $s$ by a linear
equation on the correct interval. A rational common diagonalization,
when supplied, gives rational scalar allocation data; arbitrary rational
commuting matrices need not have a rational common eigenbasis.

\begin{example}[The covariance benchmark loses order $n\log n$]\label{ex:covariance-gap}
For $f(x)=\frac12\sum_i x_i^2$ on $[0,1]^n$ at unit tolerance,
$E(P)=\tr(P^2)$. The arithmetic--geometric mean inequality gives
$D_*=n^{-n/2}$, attained by $n^{-1/2}I$, so
$\Phi=(n/4)\log_2n$. A parity contact has diameter at most $\sqrt8$,
by its midpoint gap $\|x-y\|^2/8\le1$. Its volume is at most
$\omega_n8^{n/2}$, which yields
\[
 \pconv\ge\tfrac n2\log_2n-\tfrac n2\log_2(16\pi e).
\]
Conversely, put $L=\max\{0,\lceil\frac12\log_2(n/8)\rceil\}$.
Shared coordinate squares with residual width $2^{-L}$ have total
error at most $n2^{-2L}/8\le1$ and use $nL$ binaries. Hence both
minima are $(n/2)\log_2n+O(n)$ as dimension grows. The order
$O(n\log(n+1))$ in \eqref{eq:finite-law} cannot be improved to $O(n)$
for this particular benchmark. This is an unconditional geometric
limitation, separate from computational hardness.
\end{example}

\section{Rational quadratic constructions}\label{sec:rational-quadratic}

We first make the optimal rank coefficient effective, then give an
algorithm with a dimension-only finite-accuracy loss. These are different
guarantees: the former allows an additive polynomial in input height;
the latter has an additive $O(n\log(n+1))$ independent of that height.

\subsection{Uniform rational bounds at the noncommutative-rank rate}

\begin{theorem}\label{thm:rational-ncrank}
Equation \eqref{eq:rational-rate-preview}, including its construction-time,
coefficient-length and formulation-size assertions, holds for rational
quadratic data on a rational full-dimensional box. Its parameter $r$
is the noncommutative rank, not the common nonlinear input rank used
later in \cref{thm:input-quotient}.
\end{theorem}
\begin{proof}
Let $s$ be the total binary input length. Exact elimination computes a
basis for the Hessian span. The rational algorithm of
\citet[Theorem 1.5 and Lemma 5.3]{iqs2018} returns a maximal shrunk
subspace $U$ of polynomial bit length with the same deficiency as over
$\C$. Set $V=\sum_jH_jU$ and construct rational bases for
$Z=U\cap V^\perp$, $W=V\cap U^\perp$, and $Q=(Z+W)^\perp$.
All operations are rational nullspaces and intersections. Their chosen
bases need not be orthonormal: the spaces themselves are orthogonal.
Thus $T=[Z_b\ W_b\ Q_b]$ is rational and invertible, and
\cref{lem:symmetric-shrink} gives
\[
 T^TH_jT=\begin{pmatrix}0&B_j&0\\ B_j^T&C_j&D_j\\0&D_j^T&E_j\end{pmatrix},
 \qquad 2\dim W+\dim Q=r.
\]
If $r=0$, all Hessians vanish and the affine graph needs no integers.
If $r=n$, choosing $U=V=0$ is sufficient. Rational bases, their inverses,
and the displayed coefficients have polynomial height: clear denominators
and bound the determinants of the square subsystems used by elimination.
Only a fixed number of such linear-algebra operations follows the
imported polynomial-bit algorithm.

For $v=T^{-1}x$, compute each coordinate's exact minimum and maximum on
the original box by summing the smaller or larger endpoint contribution
of each entry of $T^{-1}$. Every width is positive. Normalize this enclosing
box to $y\in[0,1]^n$, obtaining $x=c+Ay$. Retain all original bounds
on $c+Ay$ in the formulation. Translation and positive diagonal scaling
preserve the zero blocks. Write the transformed quadratics as an exact
affine part plus $\sum_{i\le t}c_{j,it}y_iy_t$, and set
\[
 C=\max_j\sum_{i\le t}|c_{j,it}|>0,\qquad
 b=\max\{0,\lceil\log_2(C/4)\rceil\},\qquad K=k+b.
\]
The integer $b$ is computed by rational comparisons with powers of two
and satisfies $b\le\operatorname{poly}(s)$. Give $Z$-coordinates depth
zero, $W$-coordinates depth $K$, and $Q$-coordinates depth
$\lceil K/2\rceil$. Every nonzero monomial has a $W$ endpoint or two
$Q$ endpoints, so its residual-width product is at most $2^{-K}$.
The shared prefix construction has output error at most
$C2^{-K}/4\le2^{-k}$ and binary count
\[
 \dim(W)K+\dim(Q)\lceil K/2\rceil
 \le\tfrac r2 k+\tfrac r2 b+\tfrac n2.
\]
All transformed fixed coefficients have polynomial bit length; the
new dyadics have $O(k+b+1)$ bits. There are $O(n^2)$ monomials, and a
monomial needs $O(L_i+L_t+1)$ rows and auxiliaries. Consequently individual
coefficients have length $\operatorname{poly}(s)+O(k)$, while the numbers
of rows and variables are
$\operatorname{poly}(s)(k+\operatorname{poly}(s)+1)$. This proves the
algorithmic upper bound, with no algebraic-number coefficient model.

For the lower bound let $D_H$ be a common denominator of the Hessian
entries and put $E_B=\prod_i\operatorname{den}(l_i)\operatorname{den}(u_i)$
for the rational box endpoints. Both logarithms are at most $s$ after
using the product of the input denominators. By \cref{lem:principal}
there is a principal coordinate set $I$ of size $r$ whose restricted
Hessians $G_j$ have full noncommutative rank. The matrices $D_HG_j$
are integral. The integral capacity bound
\citep[Theorem 2.18]{ggow2020} gives
\[
 \kappa=\operatorname{cap}(T_G)
 =D_H^{-2r}\operatorname{cap}(T_{D_HG})
 \ge D_H^{-2r}r^{-2r}.
\]
The power $2r$ is essential, since capacity is a determinant ratio.
Fixing the other coordinates leaves an $r$-box of volume
$V_I\ge E_B^{-1}$: each positive rational width is at least the
reciprocal product of its endpoint denominators. The finite capacity
bound \eqref{eq:ncrank-finite}, together with
$\omega_r^{2/r}\le4$, now gives
\[
 \eps\ge
 \frac{2^{-2p/r}}{16D_H(r+2)\sqrt{mr}\,E_B^{2/r}}.
\]
At $\eps=2^{-k}$ this implies
\begin{equation}\label{eq:rational-ncrank-lower-explicit}
 p\ge\tfrac r2 k-\tfrac r2\log_2D_H-\log_2E_B
       -\tfrac r2\log_2\!\bigl[16(r+2)\sqrt{mr}\bigr].
\end{equation}
The subtracted terms are $O(rs+r\log(mr+2)+s)$ and hence polynomial
in $s$. This lower bound permits arbitrary real convex lifts. If desired,
the principal set $I$ can be found by repeated coordinate-deletion rank
tests: while more than $r$ coordinates remain, an invertible principal
$r$-submatrix guarantees a deletion preserving rank $r$.
\end{proof}

An arbitrary rational $\eps\in(0,1]$ can be handled by the integer
$k=\lceil\log_2(1/\eps)\rceil$, computed by integer comparisons. Running
time includes the length of the rational tolerance input, not only its
numerical logarithm. Polynomial encoding permits exponential numerical
coefficient magnitudes and does not give polynomial-time MILP optimization.

\subsection{Certified rational matrix functions}\label{subsec:matrix-functions}

The finite construction requires matrix functions, but it does not require
well-separated eigenvalues. We give the numerical details because a
polynomial arithmetic-operation count alone would not establish a rational
formulation theorem. Rational Jacobi near-diagonalization is an established
spectral ingredient; see also \citet[Theorem 2]{delpia2026}.

\begin{lemma}[Rational spectral residual]\label{lem:rational-jacobi}
Given rational symmetric $A$ and rational $\sigma>0$, one can compute
rational matrices $V,B$ in deterministic polynomial time such that
\[
 V^TV=I,\qquad A=VBV^T,\qquad
 \|B-\diag(B)\|_F\le\sigma.
\]
The bit bound is polynomial in the input lengths, including the precision
length of $\sigma$. No spectral-gap hypothesis is needed.
\end{lemma}
\begin{proof}
For order one, take $V=I$. Otherwise let $N=n(n-1)$ and choose a rational
$M\ge\max(1,\|A\|_F)$, for example the larger of one and its entrywise
absolute sum. Maintain exact rational identities $A=VBV^T$ and $V^TV=I$,
starting at $B=A,V=I$. If the off-diagonal Frobenius norm $s$ exceeds
$\sigma$, select a largest off-diagonal entry $B_{ik}$. An exact
Jacobi rotation $C_*$ annihilates it and changes the squared off-norm
from $s^2$ to $s^2-2B_{ik}^2\le(1-2/N)s^2$.

Choose an exactly orthogonal rational plane rotation $C$ with
$\|C-C_*\|_2\le\sigma/(8NM)$. This is possible with polynomially many
bits: take a rational approximation of the half-angle tangent $t$ and
use $c=(1-t^2)/(1+t^2)$ and $s_\theta=2t/(1+t^2)$. The annihilating
angle may be chosen in $[-\pi/4,\pi/4]$; its tangent follows from the
two-by-two quadratic formula, and the half-angle conversion has denominator
bounded away from zero. While a rotation is needed,
$|B_{ik}|\ge\sigma/\sqrt N$, so divisions used to find that angle require
only polynomial precision depth. Rational bisection computes the square
roots and the parameter approximation.

Set $B\leftarrow C^TBC$ and $V\leftarrow VC$. Because
$\|B\|_F=\|A\|_F\le M$, the perturbation from exact conjugation is at
most $2M\|C-C_*\|_2$. Hence the new off-norm is at most
\[
 (1-1/N)s+\sigma/(4N)\le(1-3/(4N))s \quad(s>\sigma).
\]
Thus $O(N(1+\log^+(M/\sigma)))$ rotations suffice; termination is tested
by comparing rational squared norms. A common denominator for each
rotation has polynomial bit length. Each exact update adds only twice
that length to a common denominator for $B$, and once to one for $V$.
There are polynomially many updates. Norm preservation bounds the
numerators once these denominators are fixed. This proves polynomial
Turing complexity, as well as exact reconstruction and orthogonality.
\end{proof}

For positive definite $A$ with spectrum in $[a,b]$, put
$A_0=V\diag(B)V^T$. Its eigenvalues also lie in $[a,b]$, because the
entries $B_{ii}$ are Rayleigh quotients. With $\|A-A_0\|_2\le\sigma$,
\begin{equation}\label{eq:matrix-function-bounds}
 \begin{split}
 \|\log A-\log A_0\|_2&\le\sigma/a,\\
 \|A^{1/2}-A_0^{1/2}\|_2&\le\sigma/(2\sqrt a),\\
 \|A^{-1/2}-A_0^{-1/2}\|_2&\le\sigma/(2a^{3/2}).
 \end{split}
\end{equation}
The first bound follows from the resolvent formula for logarithms and
$\int_0^\infty(a+t)^{-2}\,dt=1/a$. For the second, the difference
$X=A^{1/2}-A_0^{1/2}$ solves
$A^{1/2}X+XA_0^{1/2}=A-A_0$; its integral solution gives the bound.
Multiplication by the two inverse square roots gives the third. For
symmetric arguments of norm at most $M_0$, the Duhamel formula also gives
$\|e^A-e^{A_0}\|_2\le e^{M_0}\|A-A_0\|_2$.

Thus matrix functions reduce to scalar functions of rational diagonal
entries, without approximating any particular exact eigenbasis. Scalar
square roots use rational bisection. For logarithms, dyadic range reduction
to $[1,2]$ makes the arctanh series parameter at most $1/3$; the same
series computes $\log2$. Exponentials use range reduction and a Taylor
series. When $|\log a|,|\log b|,M_0$ are polynomially bounded, prescribed
inverse-exponential-polynomial absolute errors require polynomial work
and output length. The largest diagonal entry also gives a rational
unit vector $v$, a column of $V$, with
\begin{equation}\label{eq:rayleigh-residual}
 0\le\lambda_{\max}(A)-v^TAv\le\sigma.
\end{equation}
This avoids selecting an eigenvector at a repeated largest eigenvalue.

\subsection{A polynomial-time finite-accuracy construction}

\begin{theorem}\label{thm:rational-finite}
For rational quadratic coefficients and positive rational component
tolerances, a deterministic polynomial-time algorithm constructs a rational
MILP graph relaxation with
\begin{equation}\label{eq:rational-finite-count}
 p_{\rm out}\le\pconv+O(n\log(n+1)).
\end{equation}
Its total encoding length is polynomial in the complete input length.
The implicit constant is independent of coefficients, tolerances and
output dimension. The running-time assertion is a worst-case complexity
bound, without a practical iteration-count claim.
\end{theorem}
\begin{proof}
All logarithms in this proof are natural. On the positive definite cone
use the affine-invariant metric and its distance
\[
 \langle A,B\rangle_P=\tr(P^{-1}AP^{-1}B),\qquad
 d(P,Q)=\|\log(P^{-1/2}QP^{-1/2})\|_F.
\]
This is a Hadamard manifold with sectional curvature in $[-1/2,0]$;
we use the weaker lower bound $-1$ \citep[Proposition I.1]{criscitiello2021}.
The comparison inequality of \citet[Corollary 8]{zhang2016} is the only
geometric convergence import. The following error accounting applies
it to rational stored iterates.

\emph{An exact penalty.} Omit zero Hessians from logarithmic expressions
and define
\begin{equation}\label{eq:exact-covariance-penalty}
 h(P)=\max\left\{0,\log\lambda_{\max}(P),
          \max_{j:H_j\ne0}\tfrac12\log\frac{E_j(P)}{\eps_j^2}\right\},
 \qquad F(P)=-\log\det P+nh(P).
\end{equation}
The covariance $\widehat P=e^{-h(P)}P$ is feasible and satisfies
$F(P)=-\log\det\widehat P$. Hence $\min F=-\log D_*$, attained by any
optimal feasible covariance. Each component of $h$ is geodesically convex:
use \eqref{eq:energy-geodesic}, and express the largest-eigenvalue term
as the supremum over fixed nonzero $v$ of
$\log[(v^TPv)/(v^Tv)]$. Each such Rayleigh branch is log-sum-exp along a
geodesic. In tangent coordinates $A\mapsto P^{-1/2}AP^{-1/2}$, the
half-log energy gradient and Rayleigh gradient are respectively
\begin{equation}\label{eq:energy-gradients}
 \frac{M_j^2}{\tr(M_j^2)},\quad M_j=P^{1/2}H_jP^{1/2},
 \qquad \frac{P^{1/2}vv^TP^{1/2}}{v^TPv}.
\end{equation}
They are positive semidefinite of trace one and norm at most one.
Therefore $h$ is globally one-Lipschitz and $F$ is globally
$(n+\sqrt n)$-Lipschitz.

\emph{A polynomial-radius search region.} Choose a nonnegative integer
$b$ with $\delta=2^{-b}\le1$ and
$\delta\sum_{a,c}|H_{j,ac}|\le\eps_j$ for every $j$. Rational dyadic
comparisons find $b$ of polynomial magnitude. Then $\delta I$ is feasible,
so an optimal $P_*$ has $\det P_*\ge\delta^n$ and eigenvalues at most
one. It follows that
\[
 d(I,P_*)\le\sum_i-\log\lambda_i(P_*)\le nb\log2<R:=1+nb.
\]
The closed metric ball $\mathcal B_R$ is geodesically convex. Its metric
projection is
\begin{equation}\label{eq:radial-projection}
 \Pi_R(Q)=\exp(t\log Q),\qquad
 t=\frac{R}{\max\{R,\|\log Q\|_F\}}.
\end{equation}
The reverse triangle inequality gives the lower distance to the ball,
and the radial geodesic attains it. Projection is nonexpansive on a
Hadamard manifold.

\emph{A recurrence for the actual rational iterates.} Start at $P_0=I$
and store iterates in $\mathcal B_{R+1}$. Put $D=2R+2$ and $Z=1+D$.
At the current rational $P$ obtain a symmetric tangent-coordinate matrix
$G$, with $\|G\|_F\le3n$, such that
\begin{equation}\label{eq:inexact-subgradient}
 F(P)-F(P_*)\le-\tr\!\left[G\log(P^{-1/2}P_*P^{-1/2})\right]+e.
\end{equation}
Form the conceptual exact update
$Q=P^{1/2}e^{-\eta G}P^{1/2}$, project it to $\overline Q=\Pi_R(Q)$,
and store a rational symmetric $P^+$ with $d(P^+,\overline Q)\le\xi$.
For $\xi\le1$, the next iterate again lies in $\mathcal B_{R+1}$.
The comparison inequality, using its curvature factor
$D/\tanh D\le1+D$, and nonexpansive projection give
\begin{equation}\label{eq:inexact-recurrence}
 F(P)-F(P_*)\le
 \frac{d(P,P_*)^2-d(P^+,P_*)^2}{2\eta}
 +\frac{Z\eta(3n)^2}{2}+e+\frac{2D\xi+\xi^2}{2\eta}.
\end{equation}
For the last term, both the projected point and $P_*$ are in
$\mathcal B_R$, so changing the former by distance $\xi$ changes its
squared distance to $P_*$ by at most $2D\xi+\xi^2$.
Choose
\[
 \eta=[16Z(3n)^2]^{-1},\quad
 T=\lceil16D^2/\eta\rceil,\quad e\le1/32,\quad
 \xi\le\eta/[128(D+1)].
\]
Summing \eqref{eq:inexact-recurrence} over $T$ steps telescopes on the
stored iterates. The initial-distance, curvature, oracle and rounding
terms in the average gap are at most $1/32,1/32,1/32,1/64$ respectively.
Thus some stored objective has gap below $1/4$. Selecting the smallest
reported objective, each evaluated within $1/64$, adds at most $1/32$
and gives gap below $1/2$. No global exact trajectory is being tracked;
rounding enters only the local inequality and the telescoping distances.

\emph{Implementation of the inexact oracle.} All stored spectra lie in
$[e^{-(R+1)},e^{R+1}]$. If a symmetric approximation to a positive definite
matrix with minimum eigenvalue $a$ has operator error at most $au$,
$u\le1/2$, then relative eigenvalues lie in $[1-u,1+u]$, proving
\begin{equation}\label{eq:relative-metric-bound}
 d(A,\widetilde A)\le2\sqrt n\,u.
\end{equation}
For a nonzero rational Hessian,
$E_j(P)\ge e^{-2(R+1)}\|H_j\|_F^2$; this is the minimum singular-value
bound for two-sided multiplication by $P^{1/2}$. A nonzero input entry
has magnitude at least $2^{-s}$, where $s$ is its total input length.
Every normalization denominator therefore has an inverse of at most
exponential-polynomial magnitude.

Evaluate each branch of $h$ within $\tau=1/(512n)$. For the cap branch
use the fixed rational Rayleigh vector from \eqref{eq:rayleigh-residual},
with its logarithmic shortfall at most $\tau$, rather than a possibly
unstable exact leading eigenvector. Choose a largest reported branch.
Its exact value is within $3\tau$ of $h(P)$. Approximate that branch's
tangent gradient within $\nu/n$, where $\nu=1/[128(D+1)]$, and form
$G=-I+n\widetilde\nabla h$. The full gradient error is at most $\nu$,
so \eqref{eq:inexact-subgradient} holds with
$e\le3n\tau+D\nu<1/32$ and $\|G\|_F\le3n$.
The conditioning estimates and \cref{lem:rational-jacobi} give all these
approximations in polynomial bit complexity.

The exact update $Q$ has logarithmic spectral bounds at most
$R+1+3n\eta$. Matrix functions are therefore covered by
\eqref{eq:matrix-function-bounds}. To implement projection, first replace
$Q$ by its nearby rationally diagonalizable matrix. Equation
\eqref{eq:relative-metric-bound} and nonexpansiveness control this
replacement. Then \eqref{eq:radial-projection} is a scalar spectral
formula with denominator at least $R\ge1$. Round its output symmetrically
to a fresh dyadic matrix so that \eqref{eq:relative-metric-bound} gives
the prescribed $\xi$. Polynomially many precision bits suffice and
positive definiteness is preserved. Fresh rounding prevents accumulated
symbolic denominator growth in the outer iteration.

\emph{Exact repair and a rational grid.} For the selected rational $P$,
compute a rational upper approximation $\overline\kappa$ to
\[
 \kappa=\max\{1,\lambda_{\max}(P),\max_j\sqrt{E_j(P)}/\eps_j\}
\]
within factor $e^{1/(4n)}$, using certified spectral and scalar intervals.
Then $\widehat P=P/\overline\kappa$ is exactly feasible and
\begin{equation}\label{eq:rational-covariance-quality}
 \log\det\widehat P\ge\log D_*-1.
\end{equation}
The previous radius and dyadic bounds give
$\kappa\le e^{R+1}2^b$, hence an explicit inverse-exponential-polynomial
lower bound on $\lambda_{\min}(\widehat P)$.
Apply \cref{lem:rational-jacobi} at residual at most half this lower
bound, writing $B=U^T\widehat P U$, and set
$\widetilde\lambda_i=B_{ii}/4$ and
$\widetilde P=U\diag(\widetilde\lambda_i)U^T$. Then
\begin{equation}\label{eq:rational-grid-sandwich}
 \widehat P/8\preceq\widetilde P\preceq\widehat P/2,
 \qquad \det\widetilde P\ge8^{-n}\det\widehat P.
\end{equation}
Indeed $\widetilde P$ differs from $\widehat P/4$ by at most
$\lambda_{\min}(\widehat P)/8$ in operator norm. Each energy is monotone
in positive semidefinite order: its derivative in direction $A\succeq0$
is $2\tr(H_jPH_jA)\ge0$. Thus $\widetilde P$ remains feasible.

Use the rational orthogonal matrix $U$ and rational scales
$\widetilde\lambda_i$ in \eqref{eq:covariance-grid}. Coordinate widths,
translations, transformed coefficients and prefix products are rational.
The depths are chosen by squared rational comparisons, without evaluating
an exact square root. By \eqref{eq:rational-covariance-quality} and
\eqref{eq:rational-grid-sandwich}, their sum is at most
$\Phi+B_n+3n/2+1/(2\log2)$. Combine this with
\eqref{eq:finite-law} to prove \eqref{eq:rational-finite-count}.
Finally $-\log D_*\le nb\log2$, so the depths, row counts and coefficient
lengths are polynomial in the input length.
\end{proof}

\section{Output budgets and nonlinear input coordinates}\label{sec:quadratic-bodies}

\subsection{Correlated and partially measured errors}

A finite family of positive semidefinite matrices $W_\ell\in\R^{m\times m}$
and tolerances $\eps_\ell>0$ defines the requirement
\begin{equation}\label{eq:grouped-errors}
 e^TW_\ell e\le\eps_\ell^2\quad(1\le\ell\le g),
 \qquad e=w-F(x).
\end{equation}
When $\bigcap_\ell\ker W_\ell=\{0\}$ this is an error body in
\cref{def:dimension}. Otherwise its intersection is an unbounded closed
symmetric convex set: errors in its common kernel are unmeasured. In this
subsection only, extend the same minima by using this closed set in the
graph-containment definition. This convention is explicit because the
compact-body hypothesis elsewhere remains in force.

\begin{theorem}[Grouped covariance law]\label{thm:grouped-covariance}
Replace the energies in \eqref{eq:finite-determinant} by
\[
 E_\ell(P)=\sum_{j,k}(W_\ell)_{jk}\tr(H_jPH_kP)
\]
and their bounds by $\eps_\ell^2$. Then \eqref{eq:finite-law} holds with
the same $A_n,B_n$ for \eqref{eq:grouped-errors}. For rational input,
including $W_\ell$, the algorithm of \cref{thm:rational-finite} has the
same dimension-only additive guarantee and polynomial rational encoding.
It does not require a rational factorization of $W_\ell$.
\end{theorem}
\begin{proof}
For the proof, factor $W_\ell=C_\ell^TC_\ell$ and define
$G_a=\sum_j(C_\ell)_{aj}H_j$. The energy is the nonnegative sum
$E_\ell(P)=\sum_a\tr(G_aPG_aP)$. Apply the fourth-moment identity
in \cref{lem:covariance} to these scalar combinations and sum. The
midpoint budget bounds the sum by $16\eps_\ell^2$, and every discarded
fourth-moment term is a square or squared mean. Therefore
$E_\ell(\Sigma)\le16\eps_\ell^2$, proving the covariance lower bound
as before, including for singular budgets.

For the shared grid let
$Z_{ik}=\delta_{ik}/\sqrt{\lambda_i\lambda_k}$ and
$M_j=\diag(\sqrt\lambda)U^TH_jU\diag(\sqrt\lambda)$. Every entry of
$Z$ is at most $1/(4n)$ in absolute value, so $\|Z\|_F\le1/4$, while
$e_j=\frac12\tr(M_jZ)$. Consequently
\begin{equation}\label{eq:shared-ellipsoidal-error}
 e^TW_\ell e
 =\tfrac14\sum_a\left\langle\sum_j(C_\ell)_{aj}M_j,Z\right\rangle_F^2
 \le\tfrac14E_\ell(P)\|Z\|_F^2\le\eps_\ell^2/64.
\end{equation}
Thus the grid controls an entire correlated error vector, with no output
count in the constants.

For computation retain the rational double sums. Each nonzero energy
is a positive-coefficient exponential sum along a geodesic. Its half-log
gradient is
\[
 \frac{N_\ell(P)}{E_\ell(P)},\qquad
 N_\ell(P)=\sum_{j,k}(W_\ell)_{jk}M_jM_k
       =\sum_a\left(\sum_j(C_\ell)_{aj}M_j\right)^2.
\]
The last identity proves symmetry, positive semidefiniteness and trace
$E_\ell(P)$, without using the factorization in the algorithm. An energy
is identically zero precisely when its rational value $E_\ell(I)$ is
zero; omit those logarithmic branches. Such measured output combinations
are affine and are still represented exactly. For every remaining group,
$E_\ell(P)\ge\lambda_{\min}(P)^2E_\ell(I)$, and the positive rational
$E_\ell(I)$ has polynomial encoding length. A feasible dyadic covariance
is found by exact comparisons $\delta^2E_\ell(I)\le\eps_\ell^2$.
All radius, gradient, matrix-function and repair estimates in
\cref{thm:rational-finite} now apply. The final common monomial grid
satisfies \eqref{eq:shared-ellipsoidal-error}.
\end{proof}

The certificate \cref{prop:covariance-certificate} also applies, replacing
$H_jPH_j$ by $\sum_{a,b}(W_j)_{ab}H_aPH_b$. For commuting Hessians the
scalar program of \cref{prop:commuting-covariance} uses the nonnegative
coefficients $h_{:i}^TW_\ell h_{:i}$ in place of $h_{ji}^2$.

\subsection{A specialized total-absolute-error benchmark}

For the error requirement $\|w-F(x)\|_1\le\eps$, define the Gram matrix
and correlation semidefinite bound
\begin{equation}\label{eq:l1-covariance}
 \Gamma(P)_{jk}=\tr(H_jPH_kP),\qquad
 S(P)=\max\{\tr(\Gamma(P)X):X\succeq0,\ \diag X=\mathbf1\}.
\end{equation}
The matrix $\Gamma(P)$ is positive semidefinite. The classical positive
semidefinite Grothendieck bound is
\begin{equation}\label{eq:psd-grothendieck}
 \max_{s\in\{-1,1\}^m}s^T\Gamma s
 \le S\le\tfrac\pi2\max_{s\in\{-1,1\}^m}s^T\Gamma s;
\end{equation}
see \citet{briet2010} for its attribution to Rietz and Nesterov.
One proof of the upper bound factors a correlation optimizer into unit
vectors and applies Gaussian hyperplane signs. Their expected sign product
is $\frac2\pi\arcsin X_{jk}$. The entrywise power series of
$\arcsin X-X$ is a nonnegative sum of positive semidefinite odd Hadamard
powers. Its inner product with $\Gamma\succeq0$ is nonnegative, proving
the ratio. The lower bound uses $X=ss^T$.

\begin{theorem}\label{thm:l1-covariance}
Let $D_1=\max\{\det P:0\preceq P\preceq I,S(P)\le\eps^2\}$ and
$\Phi_1=-\frac12\log_2D_1$. Then
\[
 \max\{0,\Phi_1-A'_n\}\le\pconv\le\pbin\le\Phi_1+B_n,
 \quad
 A'_n=\log_2[\omega_n(n+2)^{n/2}]
              +\tfrac n2\log_2\max\{6,n/4\}.
\]
For rational input a polynomial-time rational MILP construction again
has count at most $\pconv+O(n\log(n+1))$.
\end{theorem}
\begin{proof}
Every sign combination of the quadratic midpoint difference on a parity
contact has absolute value at most $4\eps$ in the normalization
$q_j(x-y)=\frac12(x-y)^TH_j(x-y)$. Its covariance therefore satisfies
$s^T\Gamma(\Sigma)s\le16\eps^2$. Equation \eqref{eq:psd-grothendieck}
gives $S(\Sigma)\le8\pi\eps^2<36\eps^2$. Scale by $\max\{6,n/4\}$
and apply the volume proof of \cref{thm:finite-covariance}.
For the upper grid, its shared $Z$ satisfies $\|Z\|_F\le1/4$, whence
\[
 |s^Te|\le\tfrac12\left\|\sum_j s_jM_j\right\|_F\|Z\|_F
 \le\tfrac18\sqrt{s^T\Gamma(P)s}\le\eps/8.
\]
Maximization over signs is exactly the $\ell_1$ norm.

For each fixed correlation $X$, the energy
$E_X(P)=\tr(\Gamma(P)X)$ is a grouped positive semidefinite energy.
Its nonzero half-log branch is one-Lipschitz and geodesically convex.
Their maximum $\frac12\log S(P)$ has these same properties. A nearly
optimal rational feasible $X$ supplies a near-active gradient branch,
and a certified upper value of $S(P)$ supplies the final feasibility
repair in \cref{thm:rational-finite}. We construct both next.

For rational $P$, $\Gamma$ is rational. If it is nonzero, then
$\tr\Gamma\le S\le m\tr\Gamma$. Normalize $\tr\Gamma=1$. For $m\ge3$,
parameterize $X=I+\operatorname{off}(v)$ by its $d=m(m-1)/2$ independent
off-diagonal entries. The correlation body contains the radius-$1/4$
ball around zero, since $\|\operatorname{off}(v)\|_F=\sqrt2\|v\|_2$,
and lies in the radius-$m$ ball. It has a rational strong separation
oracle: exact symmetric elimination either proves $X\succeq0$ or
returns rational $a$ with $a^TXa<0$, giving a separating inequality.
A negative pivot provides such a vector; a zero diagonal with a nonzero
off-diagonal entry provides an indefinite two-dimensional block and a
rational negative vector. Rational congruences lift it back to the original
coordinates. Fraction-free elimination gives polynomial bit bounds.
The cases $m=1,2$ are an exact scalar or interval optimization.

The weak separation--optimization theorem of
\citet[Definition (5) and Theorem (3.1)]{gls1981}, with these inner and
outer balls, returns rational $v$ within distance $\rho$ of the body
whose linear objective is within $\rho$ below its optimum. The running
time is polynomial in $\log(1/\rho)$ and the other input lengths.
Writing $s=S/\tr\Gamma$ before normalization and
$a=\tr(\Gamma X(v))$ after normalization, we have
$a\ge s-\rho$ and $\lambda_{\min}(X(v))\ge-\sqrt2\rho\ge-2\rho$.
Thus
\begin{equation}\label{eq:correlation-repair}
 X_f=\frac{X(v)+2\rho I}{1+2\rho}
\end{equation}
is rational, exactly positive semidefinite, and has unit diagonal.
Its normalized objective $a_f=(a+2\rho)/(1+2\rho)$ obeys
$0\le s-a_f\le(2m+1)\rho$. Given $0<\nu\le1/2$, choose
$\rho=\nu/[4(m+1)]$. Undoing normalization, the rational numbers
\[
 L=\tr(\Gamma(P)X_f),\qquad U=L+\nu\tr\Gamma(P)
\]
satisfy $L\le S(P)\le U$, $L\ge\tr\Gamma(P)/2$, and $U/L\le1+2\nu$.
Thus the half-log branch error is at most $\nu$, while $U$ is a
certified upper bound. These certificates follow from weak optimization;
an exact dual semidefinite solution is not assumed.

Finally, $\tr\Gamma(P)\ge\lambda_{\min}(P)^2\tr\Gamma(I)$, where the
last trace is a positive rational unless all Hessians vanish. A dyadic
feasible covariance follows from
$\delta^2m\tr\Gamma(I)\le\eps^2$. The conditioning, branch-error and
rational-repair bounds of \cref{thm:rational-finite} therefore apply with
inverse-polynomial $\nu$. When every Hessian vanishes, use the exact
affine graph instead.
\end{proof}

Rationally transformed budgets $\|Ae\|_1\le\eps$, finitely many such
budgets, and their intersections with the grouped quadratic budgets are
handled by taking the maximum of their penalty branches. The additive
count is independent of the number of groups. The hardness reduction
in \cref{sec:quadratic-hardness} uses separate component tolerances;
it does not by itself establish the same hardness for a single
$\ell_1$ budget.

\subsection{General symmetric bodies via the nonlinear output image}

For a body given by an oracle, we use the following explicit input model.
The body $K\subset\R^m$ is compact, convex, symmetric about zero, and
has known positive rational radii
$r_0B_2^m\subset K\subset R_0B_2^m$. At a rational query a strong
separation oracle returns membership or a rational valid linear inequality
strictly separating the query. Its running time and output encoding length
are polynomial in the query length and a fixed encoding of the oracle.
All subsequent oracle running-time bounds include these data.

We use one classical rounding consequence under this model. In dimension
$d$, a well-bounded symmetric body $C$ admits a polynomial-time rational
ellipsoid $E_0$ satisfying
\begin{equation}\label{eq:oracle-ellipsoid}
 E_0\subset C\subset\beta_dE_0,\qquad
 \beta_d=(d+1)\sqrt d.
\end{equation}
Here $E_0=\{z:z^TWz\le1\}$ with rational $W\succ0$. For precision
about the import, \citet[Theorem B.5]{dpv2011}, a GLS rounding theorem,
returns rational $A\succ0$ and a possibly real center $t$, with
$C\subset t+E(A)$ and either $\vol E(A)\le\eta$ or
$t+E(A)/\beta_d\subset C$. For a known inner radius $r$, choose
$\eta=\min\{1/2,(r/d)^d/2\}<\vol C$. The small-volume alternative is
impossible. Symmetry and averaging remove $t$ from both inclusions.
Take $W=d(d+1)^2A$. This proves \eqref{eq:oracle-ellipsoid} with rational
data, without computing the center or a matrix square root.

\begin{theorem}[General symmetric output body]\label{thm:general-quadratic-body}
Under the stated oracle assumptions and rational quadratic data on the
cube, a deterministic polynomial-time algorithm produces a rational MILP
with $p_{\rm out}\le\pconv(F,[0,1]^n;K)+O(n\log(n+1))$. The additive
constant is independent of output dimension; runtime still depends on
the complete input and oracle encoding.
\end{theorem}
\begin{proof}
Write $F(x)=a(x)+Cu(x)$, where $a$ is affine and $u$ lists the
$n(n+1)/2$ quadratic monomials. Choose a rational full-column-rank basis
$T$ for the image of $C$ and a rational factorization $C=TB$. Put
$d=\rank C$, $g=Bu$, and $K_{\rm eff}=\{z:Tz\in K\}$. If $d=0$,
use the exact affine graph. Otherwise rational elimination has polynomial
bit complexity, and for either formulation model
\begin{equation}\label{eq:output-image-equivalence}
 p(F,[0,1]^n;K)=p(g,[0,1]^n;K_{\rm eff}).
\end{equation}
One direction adjoins $w=a(x)+Tz$ to a reduced formulation. For the
other direction intersect an original formulation with this equation and
introduce $z$. Every exact graph point remains feasible, and
$T(z-g(x))\in K$ is precisely the reduced error condition. This does
not assume that the original formulation's errors already lay in
$\operatorname{im}T$.

Choose rational upper bounds $c_T\ge\|T\|_2$ and
$c_L\ge\|(T^TT)^{-1}T^T\|_2$ using entrywise sums plus one. Then
\[
 (r_0/c_T)B_2^d\subset K_{\rm eff}\subset(R_0c_L)B_2^d.
\]
The ambient separator at $Tz$ pulls back exactly to a separator at $z$.
Its normal cannot vanish if the query violates it, because $0\in K$.
Thus the effective body satisfies the same oracle model with polynomial
encoding. Apply \eqref{eq:oracle-ellipsoid} to obtain
$E_0\subset K_{\rm eff}\subset\alpha E_0$, where $\alpha=\beta_d$.

Let $D(t)$ be the grouped covariance determinant for $g$ with error
$tE_0$, and $\Phi(t)=-\frac12\log_2D(t)$. Scaling a covariance by
$1/\alpha$ proves
\begin{equation}\label{eq:body-rounding-loss}
 \Phi(1)\le\Phi(\alpha)+\tfrac n2\log_2\alpha.
\end{equation}
The grouped lower bound and monotonicity in the allowed error give
$\Phi(\alpha)-A_n\le\pconv(g;\alpha E_0)\le\pconv(F;K)$.
The rational construction for the inner ellipsoid uses at most
$\Phi(1)+B_n+O(n)$ binaries and admits only errors in $E_0$.
Restoring the output through $T$ proves the claim, since
$d\le n(n+1)/2$ and hence $n\log\alpha=O(n\log(n+1))$.
The formulation uses explicit rational bands from the common monomial
grid; it does not try to describe the oracle body itself by finitely
many linear inequalities.
\end{proof}

The proof also applies if the effective body has the specified bounded
strong oracle directly, even when unrelated ambient output directions
are unmeasured. It makes no claim for an effective body with no known
inner and outer radii. A factor $\alpha$ in output-norm rounding costs
at most $(n/2)\log_2\alpha$ in the covariance comparison; reducing to
the quadratic output image is what removes $m$ from this cost.

\subsection{Quotienting the common Hessian kernel}

The next parameter measures active input directions:
\[
 r_{\rm in}=\rank\begin{bmatrix}H_1\\ \vdots\\H_m\end{bmatrix}.
\]
It is distinct from the noncommutative rank of the Hessian span. In this
subsection abbreviate it by $r$.

\begin{theorem}[Nonlinear input quotient]\label{thm:input-quotient}
Under the rational and oracle hypotheses of
\cref{thm:general-quadratic-body}, there is a polynomial-time rational
MILP construction with
\begin{equation}\label{eq:input-rank-count}
 p_{\rm out}\le\pconv(F,[0,1]^n;K)+O(r\log(r+1)).
\end{equation}
For $r=0$, the exact affine graph needs no integers. Continuous size,
runtime and encoding depend on the full input dimensions.
\end{theorem}
\begin{proof}
Let $U\in\Q^{r\times n}$ have independent rows spanning all Hessian
rows and let $F_U=U^T(UU^T)^{-1}$. Then
$G_j=F_U^TH_jF_U$ satisfies $H_j=U^TG_jU$, because $F_UU$ is the
orthogonal projection onto their common row space. With
$c=\frac12\mathbf1$, write
\[
 z=U(x-c),\qquad F(x)=a(x)+g(z),\qquad
 g_j(z)=\tfrac12z^TG_jz,\qquad Z=U[-1/2,1/2]^n.
\]
The minima on the original cube and on $Z$ agree. To project a formulation,
subtract $a(x)$ from the output and use the affine visible map
$(x,w)\mapsto(U(x-c),w-a(x))$. Every $z\in Z$ has a preimage, so the
exact reduced graph is covered. Conversely, keep original continuous
$x$, its bounds, $z=U(x-c)$ and $w=a(x)+v$. These operations preserve
both models and add no integers, even when $a$ varies along a fiber.

The full-dimensional rational zonotope $Z$ has a polynomial strong
separation oracle. Membership is the rational LP
$Ut=z$, $-1/2\le t_i\le1/2$. For an infeasible query, a separating
normal can be obtained from the rational LP
\[
 v_i\ge\pm(U^Th)_i,\qquad h^Tz-\tfrac12\sum_i v_i\ge1.
\]
Strict separation and scaling guarantee feasibility of this second LP.
It returns the valid, violated inequality
$h^Ty\le\frac12\sum_i v_i$. Polynomial rational LP algorithms and
basic-solution determinant bounds give polynomial lengths. Explicit
radii follow by selecting nonsingular columns $U_J$: with
$c_{\rm inv}=1+\sum|(U_J^{-1})_{ik}|$ and
$R_Z=1+\sum|U_{ik}|$,
\[
 (2c_{\rm inv})^{-1}B_2^r\subset U_J[-1/2,1/2]^r
                 \subset Z\subset R_ZB_2^r.
\]

Apply oracle rounding to obtain $E(A)/\beta_r\subset Z\subset E(A)$,
with rational $A\succ0$. An exact factorization $A=LDL^T$ has rational
unit lower-triangular $L$ and positive diagonal $D$. Choose positive
dyadics $b_i$ by squared comparisons so that
$\sqrt{D_{ii}}\le b_i\le2\sqrt{D_{ii}}$, and put
$R=\diag(b_i)L^T$. Then
\[
 A\preceq R^TR\preceq4A,\qquad
 \beta_r^{-1}B_2^r\subset RZ\subset2B_2^r.
\]
All maps have polynomial rational encoding. The normalized domain
$\Omega=\frac12\mathbf1+\frac14RZ\subset[0,1]^r$ contains a ball of
radius $1/(4\beta_r)$ and hence a cube of side $1/[2r(r+1)]$. Therefore
\begin{equation}\label{eq:quotient-domain-volume}
 -\log_2\vol(\Omega)\le r\log_2[2r(r+1)].
\end{equation}
Its exact linear lift keeps the original $x$ and the equation
$y=\frac12\mathbf1+\frac14RU(x-c)$.

Transform the quadratics to $\Omega$ and reduce their nonlinear output
image as in \eqref{eq:output-image-equivalence}. Its dimension is at
most $r(r+1)/2$. Round its effective output body to an inner ellipsoid
with factor $\alpha=\beta_d$. Equation \eqref{eq:finite-domain-volume},
whose proof also holds for grouped budgets, and
\eqref{eq:body-rounding-loss} imply
\[
 \Phi(1)\le\pconv(F;K)+A_r-\log_2\vol(\Omega)
                           +\tfrac r2\log_2\alpha.
\]
Construct the rational grid on the containing $r$-cube, then impose the
exact domain lift and restore the original outputs. Its count is at
most $\Phi(1)+B_r+O(r)$. Equation \eqref{eq:quotient-domain-volume}
proves \eqref{eq:input-rank-count}. If the original rank is positive,
the reduced quadratic output image is nonzero, so the factor $\alpha$
is used only in positive dimension.
\end{proof}

\section{Positive quadratic structure and sharper finite bounds}\label{sec:psd-structure}

Positive semidefinite Hessians permit a first-moment trace estimate in
place of the squared Frobenius energy. When the coordinates split into
small independent blocks, this improves the dimension overhead. The
error bodies here are \emph{unconditional}: changing any coordinate sign
preserves the body. A convex unconditional body has the domination
property
\begin{equation}\label{eq:unconditional-domination}
 v\in K,\quad |u_i|\le|v_i|\ (\text{all }i)\quad\Longrightarrow\quad u\in K.
\end{equation}
Indeed, shrink one coordinate by a convex combination of its two sign
choices and repeat. This property allows a rational output rectangle to
certify membership in a curved body.

\subsection{A rational log-determinant allocation oracle}

We give a common optimization lemma for matrix blocks and scalar
allocations. Its scalar case will also support the nonlinear constructions
later in the paper. The objective and weak-optimization machinery are
classical \citep{vandenberghe1998,gls1981}; the explicit interior ball
and rational feasibility repair specify the oracle and bit model used here.

\begin{lemma}[Rational log-determinant allocation]\label{lem:block-logdet-oracle}
Let $P=(P_1,\ldots,P_B)$ be symmetric blocks of sizes $d_b$, and let
$N=\sum_bd_b\ge1$. Let $\mathcal L(P)$ be a rational linear map into
$\R^m$, and let $K$ be closed and convex with a polynomial rational
strong separation oracle and known rational $\rho_0>0$ satisfying
$\rho_0B_2^m\subset K$. Define
\[
 D=\max\left\{\prod_b\det P_b:0\preceq P_b\preceq I,
                                    \ \mathcal L(P)\in K\right\}.
\]
Given rational $0<\nu\le1$, a deterministic polynomial-time algorithm
returns exactly feasible rational positive definite blocks with product
of determinants at least $e^{-\nu}D$. The assertion requires neither
positivity of $\mathcal L$ nor symmetry or boundedness of $K$.
For blocks of size one it gives a positive rational allocation $p$ with
$Cp\in K$, $p\le\mathbf1$, and $\prod_ip_i\ge e^{-\nu}D$.
\end{lemma}
\begin{proof}
Use the independent upper-triangular entries as coordinates
$u\in\R^q$, $q=\sum_bd_b(d_b+1)/2$. Their block matrix embedding
satisfies $\|P(u)-P(v)\|_F\le2\|u-v\|_2$. Write
$\mathcal L(P(u))=Cu$, including the factor two on off-diagonal trace
coefficients, and put $c=1+\sum_{j,i}|C_{ji}|$. Choose
$\delta=2^{-b}\le1$ with $\delta Nc\le\rho_0/4$. The block identity
scaled by $\delta$ is feasible, so $D\ge\delta^N$. Every eigenvalue
of an optimizer is at least $\delta^N$, because all its eigenvalues
are at most one. Put
\[
 \begin{gathered}
 a=\delta^N/4,\quad B_0=N(Nb+2)+4,\quad P_0=(\delta/2)I,\\
 t_0=-N(b+1)-2,\qquad
 \sigma=\min\{\delta/(32N),\rho_0/(8c),1/4\}.
 \end{gathered}
\]
Let $u_0$ encode $P_0$. The compact convex hypograph
\begin{equation}\label{eq:logdet-hypograph}
 Q=\left\{(u,t):aI\preceq P_b(u)\preceq I,\quad Cu\in K,\quad
                -B_0\le t\le g(P):=\sum_b\log\det P_b\right\}
\end{equation}
has optimal last coordinate $\log D$. It contains the Euclidean ball
of radius $\sigma$ about $(u_0,t_0)$. To check this, coordinate
perturbations of norm $\sigma$ change block spectra by at most
$2\sigma$, leaving every eigenvalue between $\delta/4$ and one, and
above $a$. The center's image has norm at most $\rho_0/8$, and the
image perturbation is at most $c\sigma\le\rho_0/8$. On this ball,
the independent-coordinate gradient of $g$ has diagonal entries from
$P_b^{-1}$ and twice its off-diagonal entries, with norm at most
$8N/\delta$. Its variation is at most $1/4$. At the center
$g(P_0)-t_0\ge2$, whereas the last-coordinate perturbation is at most
$1/4$; also $t_0+B_0\ge3$. These prove all the ball inclusions.
An outer ball about the same center of radius $2(B_0+q+1)$ follows
from the spectral and last-coordinate bounds. All constants have
polynomial rational length.

A weak separation oracle for $Q$ checks the linear last-coordinate
bounds and the spectral caps exactly; rational symmetric elimination
supplies the latter separators as in the proof of
\cref{thm:l1-covariance}. It pulls back a violated separator for $Cu\in K$.
For a query passing these tests, approximate $g(P)$ within $\eta/2$ by
rational $\widehat g$. If $t>\widehat g+\eta/2$, the rational tangent
\begin{equation}\label{eq:logdet-tangent}
 t'\le\widehat g+\eta/2+
                  \sum_b\tr[P_b^{-1}(P_b'-P_b)]
\end{equation}
is valid and strictly separates the query. Otherwise lowering $t$ to
$\min(t,g(P))$ proves distance at most $\eta$ to $Q$. This correction
respects the lower last-coordinate bound because
$g(P)\ge N\log a=-N(Nb+2)\log2>-B_0$.
Only logarithms of exact positive rational determinants are approximated;
the inverse matrices in the tangent are exact rational matrices. The
spectral lower cap ensures polynomial precision depth. The tangent
normal has norm at least one through its $t$ coefficient; other normals
can be divided by their nonzero infinity norm. Thus this is the weak
separation convention of \citet[Theorem (3.1)]{gls1981}.

Weakly maximize the last coordinate to tolerance
$\rho=\nu\sigma/[4(B_0+\sigma+1)]$. The result is rational
$y=(u_y,t_y)$ with $\operatorname{dist}(y,Q)\le\rho$ and
$t_y\ge\log D-\rho$. Define the rational repaired point
\begin{equation}\label{eq:central-ball-repair}
 y_f=\frac{\sigma y+\rho(u_0,t_0)}{\sigma+\rho}.
\end{equation}
For some $z\in Q$, $\|y-z\|\le\rho$. Equation
\eqref{eq:central-ball-repair} is the convex combination of $z$ and
$(u_0,t_0)+(\sigma/\rho)(y-z)$ with weights
$\sigma/(\sigma+\rho)$ and $\rho/(\sigma+\rho)$. The second point
lies in the known interior ball. Thus $y_f\in Q$ exactly. Moreover
$0\le\log D-t_0\le B_0$, so its objective loss is at most
$\rho+(\rho/\sigma)B_0\le\nu$. The first coordinates give the claimed
rational blocks. The same proof with $d_b=1$ proves the scalar case.
\end{proof}

\subsection{Independent positive semidefinite blocks}

Suppose the original coordinates have a common partition of sizes
$d_1,\ldots,d_B$, with $N=\sum_bd_b$, and
\[
 H_j=\diag(H_{j1},\ldots,H_{jB}),\qquad H_{jb}\succeq0.
\]
Any block absent from all quadratic parts can first be removed after
affine output subtraction. If all blocks disappear, the graph is affine
and needs no integers; henceforth $N\ge1$. The Hessians need not commute
within a block. For positive definite block variables set
\[
 \begin{gathered}
 E(P)_j=\sum_b\tr(H_{jb}P_b),\qquad
 \Phi_K=-\tfrac12\log_2D_K,\\
 D_K=\max\left\{\prod_b\det P_b:0\preceq P_b\preceq I,\ E(P)\in K\right\}.
 \end{gathered}
\]
Write $c_b=\max\{4,d_b/4\}$ and
\begin{equation}\label{eq:block-constants}
 A_{\mathcal B}=\log_2[\omega_N(N+2)^{N/2}]
                  +\tfrac12\sum_bd_b\log_2c_b,
 \qquad B_{\mathcal B}=N+\sum_bd_b\log_2d_b.
\end{equation}

\begin{theorem}[Positive block allocation]\label{thm:block-psd}
For an unconditional error body $K$,
\begin{equation}\label{eq:block-psd-law}
 \max\{0,\Phi_K-A_{\mathcal B}\}\le\pconv\le\pbin
                                       \le\Phi_K+B_{\mathcal B}.
\end{equation}
For rational data and the strong-oracle hypotheses on $K$, a
polynomial-time rational MILP construction satisfies
\[
 p_{\rm out}\le\pconv+O\!\left(\sum_bd_b\log(d_b+1)\right).
\]
If $r_b=\rank[H_{1b};\ldots;H_{mb}]$, the sharper intrinsic bound is
\begin{equation}\label{eq:block-rank-count}
 p_{\rm out}\le\pconv+
       O\!\left(\sum_{b:r_b>0}r_b\log(r_b+1)\right).
\end{equation}
For a tolerance box the feasibility condition is simply
$\sum_b\tr(H_{jb}P_b)\le\eps_j$ for every output.
\end{theorem}
\begin{proof}
The midpoint error of a parity contact is the nonnegative vector
$J_j(x,y)=\frac18\sum_b(x_b-y_b)^TH_{jb}(x_b-y_b)\in K$.
For independent uniform contact points with full covariance $\Sigma$,
convexity gives $\mathbb EJ\in K$ and
$(\mathbb EJ)_j=\frac14\sum_b\tr(H_{jb}\Sigma_{bb})$.
Set $P_b=\Sigma_{bb}/c_b$. Each block covariance is capped by
$(d_b/4)I$ using its original cube coordinate range. Positivity gives
$0\le E(P)\le\mathbb EJ$, so \eqref{eq:unconditional-domination}
makes $P$ feasible. The block determinant inequality
$\det\Sigma\le\prod_b\det\Sigma_{bb}$ follows by iterated Schur
complements. Combining it with the covariance-volume inequality bounds
every contact by $2^{A_{\mathcal B}}\sqrt{D_K}$, proving the lower bound.

Diagonalize each $P_b$ within its own block. The corresponding rotated
coordinate has width at most $\sqrt{d_b}$. Choose residual widths
$h_i\le\sqrt{\lambda_i/d_b}$, giving at most
$-\frac12\log_2\prod_b\det P_b+B_{\mathcal B}$ binaries. The normalized
shared residual error matrix $Z_b$ has entries at most $1/(4d_b)$,
hence operator norm at most $1/4$. The scaled transformed Hessian
$M_{jb}$ is positive semidefinite with trace $\tr(H_{jb}P_b)$. Therefore
\[
 |e_j|\le\tfrac12\sum_b|\tr(M_{jb}Z_b)|
 \le\tfrac18\sum_b\tr(H_{jb}P_b)=E(P)_j/8.
\]
For $M\succeq0$, the inequality $|\tr(MZ)|\le\tr(M)\|Z\|_2$ follows
from $-\|Z\|_2I\preceq Z\preceq\|Z\|_2I$. Unconditional domination
proves the whole-body error guarantee, and exact monomial values prove
graph containment. An optimal allocation proves the finite upper bound.

For computation apply \cref{lem:block-logdet-oracle} with $\nu=1$.
It gives exactly feasible rational blocks of determinant product at
least $e^{-1}D_K$. Apply \cref{lem:rational-jacobi} separately within
each block and use \eqref{eq:rational-grid-sandwich}; the resulting
rational grid covariance loses at most a factor $8^N$ in determinant.
Its trace vector is dominated by the original trace vector, so it remains
feasible in $K$. The extra count is $O(N)$. Since the global volume
factor in \eqref{eq:block-constants} has logarithm $O(N)$, this proves
the first rational guarantee with polynomial encoding.

Finally quotient the common kernel separately inside each original
block by the construction of \cref{thm:input-quotient}. Congruence
preserves positive semidefiniteness. Because the original input blocks
are disjoint, the projected domain is a product of zonotopes. Normalize
each nonzero block quotient separately to
$\Omega_b\subset[0,1]^{r_b}$ with
$\vol(\Omega_b)\ge[2r_b(r_b+1)]^{-r_b}$. The lower contact proof now
covers the product of these volumes; its additional loss is at most
$\sum_br_b\log_2[2r_b(r_b+1)]$. The upper grid is built on the
containing product cube and restricted by the original continuous input
variables. Applying the preceding finite bounds with dimensions $r_b$
proves \eqref{eq:block-rank-count}. Zero-rank blocks remain only in
affine output terms and cost no integers.
\end{proof}

For separate component tolerances there is also a direct product-cone
version of the geodesic algorithm. With
$E_j(P)=\sum_b\tr(H_{jb}P_b)$ use
\[
 h(P)=\max\{0,\max_b\log\lambda_{\max}(P_b),
                    \max_{j:E_j(I)>0}\log(E_j(P)/\eps_j)\},
 \quad F(P)=-\sum_b\log\det P_b+Nh(P).
\]
The repair is $P_b\mapsto e^{-h(P)}P_b$. Along a block geodesic the
nonzero trace energies are positive-coefficient exponential sums, and
their normalized gradient blocks are
$P_b^{1/2}H_{jb}P_b^{1/2}/E_j(P)$ with total trace one. Thus the same
$2N$ Lipschitz bound applies. A feasible dyadic allocation has
$\delta E_j(I)\le\eps_j$, so the optimizer lies in a polynomial-radius
product-cone ball, and
$E_j(P)\ge\min_b\lambda_{\min}(P_b)E_j(I)$ controls denominators.
Block diagonal matrices form a totally geodesic subspace; logarithms,
exponentials and identity-centered radial projection preserve it.
Consequently the full inexact proof of \cref{thm:rational-finite} applies
with $N$ replacing $n$, performing all numerical operations within blocks.
This supplies the same rational allocation guarantee. The Euclidean
hypograph oracle above additionally handles arbitrary unconditional
strong-oracle budgets directly.

\subsection{Diagonal Hessians and an explicit linear overhead}

For original-axis diagonal positive semidefinite Hessians, the case with
no active coordinate is the exact affine graph. Otherwise retain only
$r\ge1$ active coordinates and write
$f_j(x)=\frac12\sum_i h_{ji}x_i^2+\text{affine}$, with $h_{ji}\ge0$.
Define
\[
 D_{\rm tr}=\max\left\{\prod_ip_i:0\le p_i\le1,
                                    \ Hp\in K\right\},
 \quad\Phi_{\rm tr}=-\tfrac12\log_2D_{\rm tr},\quad
 A_r^{\rm tr}=r+\log_2[\omega_r(r+2)^{r/2}].
\]
Here $H=(h_{ji})$ is the coefficient array rather than a Hessian matrix.
For a tolerance box, the budget is $\sum_i h_{ji}p_i\le\eps_j$.

\begin{corollary}[Diagonal trace law]\label{cor:diagonal-psd}
For any unconditional error body,
\begin{equation}\label{eq:diagonal-law}
 \max\{0,\Phi_{\rm tr}-A_r^{\rm tr}\}\le\pconv\le\pbin
 \le\Phi_{\rm tr}+r,\qquad A_r^{\rm tr}<4r.
\end{equation}
Under rational strong-oracle input, a polynomial-time rational MILP has
\begin{equation}\label{eq:diagonal-count}
 p_{\rm out}\le\pconv+5r+1.
\end{equation}
\end{corollary}
\begin{proof}
Use blocks of size one in \cref{thm:block-psd}: the contact allocation
is $p_i=\Sigma_{ii}/4$, and Hadamard's inequality gives the lower
constant displayed. Alternatively, this is the direct first-moment
identity $\mathbb EJ=H\diag(\Sigma)/4\in K$.
A feasible allocation has coordinate depths
$L_i=\lceil\frac12\log_2(1/p_i)\rceil$. Its shared square error is
componentwise bounded by $Hp/8$, proving the upper bound. The Gaussian
volume estimate gives
\[
 A_r^{\rm tr}\le r+\tfrac r2\log_2[2\pi e(1+2/r)]<4r.
\]
The scalar case of \cref{lem:block-logdet-oracle} provides rational $p$
with product at least $e^{-1}D_{\rm tr}$, adding at most
$1/(2\log2)<1$ to the real-valued count bound. No Jacobi approximation
is needed in the original axes. This proves \eqref{eq:diagonal-count}.
\end{proof}

For tolerance boxes an elementary log-coordinate algorithm avoids a
matrix or body oracle. Choose $\delta=2^{-b}$ with
$\delta\sum_i h_{ji}\le\eps_j$, and put $R=rb+1$. Every optimal
allocation has $p_i\ge\delta^r$, so its logarithms lie in $[-R,0]^r$.
Define
\[
 h(u)=\max\{0,\max_i u_i,
     \max_{j:\sum_i h_{ji}>0}\log(\sum_i h_{ji}e^{u_i}/\eps_j)\},
 \qquad F(u)=-\sum_i u_i+rh(u).
\]
The repair $p_i=e^{u_i-h(u)}$ is feasible and has negative log product
$F(u)$. Its minimum on the box equals $-\log D_{\rm tr}$; its
subgradients have norm at most $2r$ because each row-branch gradient
is a probability vector. Set $D_0=rR$, $G_0=3r$,
$\eta=(16G_0^2)^{-1}$ and $T=\lceil16D_0^2/\eta\rceil$. Starting
at zero, approximate branch values within $1/(128r)$ and the full
selected subgradient within $1/(64D_0)$, then take a rational step and
clip exactly to $[-R,0]^r$. The approximate subgradient inequality has
error at most $1/32$. The usual squared-distance expansion for projected
Euclidean subgradients telescopes to mean objective gap at most
$1/32+1/32+1/32=3/32$. Choosing the least objective estimate at accuracy
$1/64$ gives gap at most $1/8$.
All nonzero row sums on the box are at least
$e^{-R}\sum_i h_{ji}$, so scalar evaluation has polynomial precision
depth. Store gradients on a fixed sufficiently fine dyadic grid; the
polynomial number of rational steps has polynomial encoding length.
Choose rational $q_i>0$ with
$|\log q_i-u_i|\le1/(32r^2)$ and set
\[
 \kappa=\max\{1,\max_iq_i,\max_j(\sum_i h_{ji}q_i)/\eps_j\},
 \qquad p_i=q_i/\kappa.
\]
This is exactly feasible, and the $2r$ Lipschitz bound gives
$-\sum_i\log p_i=F(\log q)\le-\log D_{\rm tr}+1$.
Thus this scalar procedure proves the same rational count independently
of general weak optimization.

The trace benchmark differs from the Frobenius benchmark. In
\cref{ex:covariance-gap}, $D_{\rm tr}=r^{-r}$ and
$\Phi_{\rm tr}=(r/2)\log_2r$, which agrees with the true count up to
$O(r)$. Also, diagonal structure here means the original box axes, or
a signed permutation of them. A simultaneous diagonalization in an
arbitrary rotated basis can produce a correlated domain and does not
by itself prove \eqref{eq:diagonal-count}.

\subsection{Independent integer features and forest differences}

An explicit domain-volume estimate can preserve a linear overhead even
when the square features overlap in the original coordinates.
Let $T\in\Z^{r\times n}$ have full row rank, with every row active in
\[
 f_j(x)=\tfrac12\sum_{i=1}^r a_{ji}(T_ix)^2+\ell_j^Tx+b_j,
 \qquad a_{ji}\ge0,\quad x\in[0,1]^n.
\]
Put $w_i=\sum_k|T_{ik}|\ge1$, $l_i=\sum_{k:T_{ik}<0}T_{ik}$,
$C_{ji}=a_{ji}w_i^2$, and
\[
 \begin{gathered}
 \Omega=\{\diag(w)^{-1}(Tx-l):x\in[0,1]^n\}\subset[0,1]^r,
 \qquad V=\vol_r(\Omega),\\
 D=\max\{\prod_ip_i:0\le p_i\le1,\ Cp\in K\},
 \qquad\Phi=-\tfrac12\log_2D.
 \end{gathered}
\]

\begin{theorem}[Integer-feature precision]\label{thm:integer-features}
For an unconditional error body $K$ and $r\ge1$,
\begin{equation}\label{eq:feature-finite}
 \max\{0,\Phi-A_r^{\rm tr}+\log_2V\}\le\pconv\le\pbin\le\Phi+r.
\end{equation}
For rational coefficients and the strong-oracle assumptions on $K$, a
polynomial-time rational MILP satisfies
\begin{equation}\label{eq:feature-count}
 \begin{split}
 p_{\rm out}&\le\pconv+A_r^{\rm tr}+r-\log_2V+1/(2\log2)\\
 &\le\pconv+5r+\sum_i\log_2w_i+1.
 \end{split}
\end{equation}
The common nonlinear input rank is $r$. The feature representation is
part of the input; discovering such a positive decomposition is not
part of the theorem.
\end{theorem}
\begin{proof}
Choose a nonsingular $r$-column minor $T_I$. Fixing the other input
coordinates at zero gives a parallelotope contained in $\Omega$, so
\begin{equation}\label{eq:integer-feature-volume}
 V\ge\frac{|\det T_I|}{\prod_iw_i}\ge\frac1{\prod_iw_i}.
\end{equation}
The last step uses integrality of its nonzero determinant. A minor can
be found by exact elimination, and its actual determinant gives a sharper
bound if useful; no volume computation for the whole zonotope is needed.
In coordinates $u_i=(T_ix-l_i)/w_i$, the quadratic parts are
$\frac12\sum_i a_{ji}(w_iu_i+l_i)^2$ with diagonal Hessian entries
$C_{ji}$. Subtracting the original affine output before projecting and
restoring it afterward gives exact equality of formulation minima,
by the fiber argument in \cref{thm:input-quotient}.

The proof of \cref{cor:diagonal-psd} bounds each parity contact inside
$\Omega$ by $2^{A_r^{\rm tr}}\sqrt D$. Their cover has volume $V$,
which proves the lower bound. The square-prefix construction on the
containing cube uses at most $\Phi+r$ binaries, with error dominated
by $Cp/8$. Keep original continuous $x$, its box bounds and the equation
for $u$ to restrict to the actual domain. This proves the upper bound
and preserves exact graph points. A rational allocation with product at
least $e^{-1}D$ from \cref{lem:block-logdet-oracle} adds less than one
to the count. Equation \eqref{eq:integer-feature-volume} proves the
second inequality in \eqref{eq:feature-count}.
Finally, the common kernel of the Hessians is $\ker T$: positivity
and activity force each $T_ix=0$ when all quadratic forms vanish.
Thus its codimension is $r$.
\end{proof}

Rational feature rows may be cleared of denominators, with their row
scales absorbed into $a_{ji}$, then divided by their greatest common
divisors to make them primitive. These are polynomial-bit operations.
The resulting $w_i$ may be large, so the theorem does not imply a
uniform linear overhead for arbitrary rational changes of coordinates.
With a uniform row-length bound $w_i\le s$, it gives overhead at most
$(5+\log_2s)r+1$, allowing overlapping features and noncommuting Hessians.

For a forest, the row-length estimate has a sharper exact-volume form.
Orient a forest on $n$ vertices and let $B$ be its $r$-row edge incidence
matrix, after removing edges absent from every output. If no edge remains,
the graph is affine and has an exact LP; assume $r\ge1$ below. For
$f_j(x)=\frac12\sum_ea_{je}(x_u-x_v)^2+\text{affine}$, every incidence
row has width two. Let $n_c$ be the component vertex counts, including
isolated vertices. Then the normalized incidence domain satisfies
\begin{equation}\label{eq:forest-volume}
 V_F=\vol_r\bigl((B[0,1]^n+\mathbf1)/2\bigr)
                    =2^{-r}\prod_cn_c\ge2^{-r}.
\end{equation}
For a tree on $s$ vertices, subtracting the minimum coordinate gives
the unique cube representative of its incidence image with minimum zero.
Thus the image is tiled by the images of its $s$ faces $x_v=0$.
Each has volume $|\det B_{-v}|=1$, by repeated leaf-column expansion.
Distinct face images overlap only where two coordinates vanish, a set
of dimension at most $s-2$. The image volume is therefore $s$.
Different components give a product, and scaling all $r$ coordinates
by one half proves \eqref{eq:forest-volume}. This also proves incidence
row independence; positivity and activity make the common nonlinear
rank $r$.

\begin{corollary}[Forest precision]\label{cor:forest-precision}
In the preceding forest setting, take $C_{je}=4a_{je}$ in the feature
allocation. Equation \eqref{eq:feature-finite} holds with the exact
volume $V_F$, and rational oracle input admits
\[
 p_{\rm out}\le\pconv+A_r^{\rm tr}+r-\log_2V_F+1/(2\log2)
                      \le\pconv+6r+1.
\]
\end{corollary}
A path or star may have one original-coordinate Hessian block of
arbitrarily large rank. Thus the forest conclusion is distinct from
the fixed-block-rank guarantee. Cyclic edge systems need not supply
independent diagonal square coordinates, so the forest proof does not
extend by simply choosing a subset of their edges.

\subsection{Why the actual domain must be preserved}

\begin{example}[A thin domain defeats independent-block transfer]\label{ex:thin-domain}
Let $L\ge1$, $\eps=4^{-L}$, and $\delta=\eps/16$. On the original
unit square consider $f_1(x)=x_1^2$ and
$f_2(x)=(x_1+\delta x_2)^2$ with separate tolerance $\eps$. The coordinates
$z_1=x_1$, $z_2=(x_1+\delta x_2)/(1+\delta)$ map the square to a
parallelogram $\Omega\subset[0,1]^2$ of area $\delta/(1+\delta)$.
The outputs become the independent positive squares
$g_1(z)=z_1^2$, $g_2(z)=(1+\delta)^2z_2^2$.

On the actual domain, $L$ binaries suffice: a shared square-prefix
approximation $v$ of $x_1^2$ has error at most $\eps/4$. Introduce
continuous $0\le d\le2\delta+\delta^2$ and set $w_1=v,w_2=v+d$.
Every exact graph point is retained by taking
$d=2\delta x_1x_2+\delta^2x_2^2$, and the second output error is at
most $\eps/4+2\delta+\delta^2<\eps$. Conversely, on the entire product
square, each parity contact has coordinate widths at most
$2\sqrt\eps$ and $2\sqrt\eps/(1+\delta)$. Its area is at most
$4\eps/(1+\delta)$, so
\[
 \pconv(g,[0,1]^2)\ge\log_2\frac{1+\delta}{4\eps}>2L-2.
\]
Replacing the actual domain by its enclosing product therefore creates
an unbounded count gap in fixed dimension. Original disjoint blocks,
the product-zonotope argument, and the domain-volume terms above are
substantive hypotheses. Copy variables do not supply independent domains.
For $\delta=1/M$, the primitive integer features are $(1,0)$ and
$(M,1)$, whose row-length term $\log_2(M+1)$ is consistent with this
obstruction.
\end{example}

\section{Limits on approximating the minimum integer count}\label{sec:quadratic-hardness}

The preceding algorithms construct a formulation with a near-minimal
integer count. This is separate from optimizing a nonlinear objective
or solving the resulting MILP. Even approximating the minimum count
itself has a computational barrier.

For an unweighted graph $G=(V,E)$, let
\[
 f_G(x)=\sum_{\{i,j\}\in E}(x_i-x_j)^2,\qquad x\in[0,1]^n,
\]
and let $M(G)$ be its maximum cut size. Successively maximizing over
one cube coordinate moves it to an endpoint without decreasing a convex
function, so $\max f_G=M(G)$. Also $f_G(1-x)=f_G(x)$ and
$f_G(\frac12\mathbf1)=0$.

\begin{lemma}[Zero dimension]\label{lem:zero-count-maxcut}
At positive scalar tolerance $\eps$,
\[
 \pconv(f_G)=0\quad\Longleftrightarrow\quad\pbin(f_G)=0
               \quad\Longleftrightarrow\quad M(G)\le\eps.
\]
\end{lemma}
\begin{proof}
If $M(G)\le\eps$, the LP $x\in[0,1]^n$, $0\le w\le\eps$ contains
the graph and has error at most $\eps$. Conversely, a zero-integer
convex lift contains the midpoint of witnesses for the maximum-cut
vertices $v$ and $1-v$. Its visible point is
$(\frac12\mathbf1,M(G))$, with error $M(G)$.
\end{proof}

Taking $\eps=k-1/2$ for an integer cut threshold $k\ge1$ makes zero
dimension equivalent to $M(G)<k$. Thus zero recognition is coNP-complete
on this restricted graph-quadratic family, using the established
unweighted Max-Cut completeness theorem \citep{garey1976}. This is not
a coNP-membership assertion for arbitrary convex-lift minimization.
Any polynomial-time valid-formulation algorithm with a finite
multiplicative approximation guarantee on all instances would distinguish
zero from positive optimum, and hence would imply $P=NP$.

\begin{theorem}[Positive-optimum approximation hardness]\label{thm:count-hardness}
For every fixed $\delta>0$, unless $P=NP$, no polynomial-time algorithm
always constructs a valid graph relaxation with count at most
$p_{\min}+O(N^{1-\delta})$, where $N$ is the input dimension and
$p_{\min}$ is either $\pconv$ or $\pbin$. Even under the promise
$p_{\min}\ge1$, neither that additive guarantee nor a multiplicative
factor $O(N^{1-\delta})$ is possible. These conclusions hold for
rational convex quadratic outputs on disjoint input blocks, with every
component tolerance fixed at one.
\end{theorem}
\begin{proof}
Replicate the graph quadratic on $t$ independent input blocks and give
each output tolerance $k-1/2$. If $M(G)<k$, the product of the zero-integer
LPs in \cref{lem:zero-count-maxcut} has optimum zero. If $M(G)\ge k$,
choose a maximum-cut vertex $v$ and independently select $v$ or $1-v$
in each block. This gives $2^t$ exact graph points, any two of which
have an inadmissible graph midpoint in some output. Any integer lift
therefore needs distinct parities for these points, forcing $p\ge t$.
The argument uses no boundedness of integer ranges or continuous lift
size. Dividing each graph output by $k-1/2$ changes its tolerance to
one and preserves the two cases, with polynomial rational encoding.

To obtain a positive optimum, append $z\in[0,1]$ and the output $8z^2$
with unit tolerance. This scalar graph has minimum exactly one: its
endpoint midpoint has error two, while one binary $\beta$ suffices by
\[
 z=(\beta+r)/2,\quad 0\le r\le1,\quad v=\beta r,
 \quad s\ge0,\quad s\ge2r-1,\quad s\le r,
 \quad w=2(\beta+2v+s).
\]
Represent $v=\beta r$ by its exact linear binary-product inequalities.
The residual square triangle has error at most $1/4$, so this output
has error at most $1/2$ and contains the exact graph. In the low-Max-Cut
case the full optimum is one. In the high case the extra scalar endpoints
double the incompatible packing, forcing $p\ge t+1$.

It suffices to consider $0<\delta\le1$. If an additive guarantee has
constant $C$, choose a fixed integer $q>1/\delta$ and $t=n^q$.
For all sufficiently large $n$,
$C(tn+1)^{1-\delta}<t$. Thus in the low case an additive algorithm
returns fewer than $t+1$ integers, and so does a multiplicative algorithm
with that factor. In the high case every valid result uses at least
$t+1$. This decides Max-Cut in polynomial time, since replication is
a fixed polynomial in $n$. Finitely many smaller dimensions may be
handled by enumeration or isolated-vertex padding. Larger $\delta$
would imply an already excluded guarantee. The same reasoning without
the scalar appendage proves the unrestricted additive assertion.
\end{proof}

Since the common nonlinear input rank is at most $N$, replacing $N$
by that rank in any of these power-sublinear guarantees would also
contradict \cref{thm:count-hardness}. This does not exclude, for example,
an additive $N/\log N$ algorithm and does not prove a sharp linear
barrier. Together with \cref{thm:input-quotient}, it locates the attainable
dimension exponent up to subpolynomial factors. The separate geometric
example \cref{ex:covariance-gap} explains the logarithm lost by the
particular Frobenius benchmark, not a universal impossibility for every
other construction or invariant.
